% ৬. ডেটা টাইপ
\section{ডেটা টাইপ}
ডেটা টাইপ নির্ধারণ করে একটি ফিল্ডে কী ধরনের মান রাখা যাবে। সঠিক ডেটা টাইপ ব্যবহার করলে স্টোরেজ সাশ্রয় হয়, ভুল ডেটা এন্ট্রি কমে এবং কুয়েরি/গণনা সহজ হয়।

\subsection{প্রয়োগ}
\begin{center}
\begin{tabular}{|p{0.22\textwidth}|p{0.46\textwidth}|p{0.20\textwidth}|}
\hline
ডেটা টাইপ & ব্যবহার & উদাহরণ \\
\hline
Short Text & ছোট লেখা, কোড, মোবাইল নম্বর & Name, Mobile \\
\hline
Long Text/Memo & বড় বিবরণ বা মন্তব্য & Address, Note \\
\hline
Number & গাণিতিক হিসাবযোগ্য সংখ্যা & Marks, Quantity \\
\hline
Date/Time & তারিখ ও সময় & BirthDate \\
\hline
Currency & টাকা-পয়সা & Salary, Price \\
\hline
AutoNumber & স্বয়ংক্রিয় অনন্য নম্বর & StudentID \\
\hline
Yes/No & সত্য/মিথ্যা বা হ্যাঁ/না মান & Paid, Active \\
\hline
Hyperlink & ওয়েব/ইমেইল লিংক & Website \\
\hline
\end{tabular}
\end{center}

\begin{keypointbox}{ডেটা টাইপ নির্বাচন}
যে ফিল্ডে যোগ, বিয়োগ বা গড় বের করতে হবে সেটি Number/Currency হওয়া উচিত। যে নম্বরের শুরুতে 0 থাকতে পারে বা গণনা দরকার নেই, সেটি Text রাখা ভালো।
\end{keypointbox}

\begin{boardbox}{বোর্ড ফোকাস}
প্রশ্নে যদি ``মোবাইল নম্বরের জন্য কোন ডেটা টাইপ?'' আসে, উত্তর হবে Text/Short Text; কারণ মোবাইল নম্বর গাণিতিক হিসাবের জন্য নয় এবং শুরুতে 0 থাকতে পারে।
\end{boardbox}
